\documentclass[11pt,a4paper]{article}

% Core packages
\usepackage[utf8]{inputenc}
\usepackage[T1]{fontenc}
\usepackage{amsmath,amssymb,amsthm}
\usepackage{mathtools}
\usepackage{longtable,booktabs,array}
\usepackage{geometry}
\geometry{margin=1in}
\usepackage{graphicx}
\usepackage[round,sort&compress]{natbib}
\usepackage{hyperref}
\usepackage{xcolor}

% TikZ and PGFPlots for figures
\usepackage{tikz}
\usepackage{pgfplots}
\pgfplotsset{compat=1.18}
\usetikzlibrary{patterns, positioning, arrows.meta, shapes.geometric, backgrounds, fit}

% Misc
\usepackage{enumitem}
\usepackage{microtype}
\usepackage{caption}

% Custom commands
\newcommand{\R}{\mathbb{R}}
\DeclareMathOperator{\Var}{Var}
\DeclareMathOperator{\Corr}{Corr}
\DeclareMathOperator{\argmax}{arg\,max}

\title{\textbf{Appendix B: Cognitive Extensions}\\[6pt]
\large Attentional Fidelity, Executive Control, Metacognition, Social Cognition, and Multimodal Transfer}
\author{Supplementary Material for ``The Σ-Model Model''}
\date{Version 3.0+ \quad$\cdot$\quad March 2026}

\begin{document}

\maketitle

% ============================================================
\section*{Overview}\label{app:overview}
% ============================================================

This appendix contains the cognitive extensions to the core Σ-Model
framework presented in the main paper. These extensions build upon
the three core state variables ($\delta_A$, $\beta_A$, $\sigma_A$)
and are not required for understanding the basic Σ-Model system.
Readers seeking the minimal Σ-Model framework need only study the
main paper through Section~3.

The extensions are organised as follows:
\begin{itemize}[nosep]
\item \textbf{Extension 1:} Attentional Fidelity $\alpha_A(d,t)$ (Section~\ref{app:attention})
\item \textbf{Extension 2:} Collective Schema Field (Social Cognition) (Section~\ref{app:social})
\item \textbf{Extension 3:} Executive Control State $\Xi_A(t)$ (Section~\ref{app:executive})
\item \textbf{Extension 4:} Self-Model of Schema Coherence $\hat{M}_A(d,t)$ (Section~\ref{app:metacognition})
\item \textbf{Extension 5:} Cross-Modal Schema Transfer (Section~\ref{app:multimodal})
\item \textbf{Extension 6:} Benchmark Validity and Reliability (Section~\ref{app:validity})
\end{itemize}

% ============================================================
\section{Extension 1: Attentional Fidelity $\alpha_A(d,t)$}\label{app:attention}
% ============================================================

\subsection{Definition}

Attentional fidelity $\alpha_A(d,t) \in [0,1]$ quantifies the degree
to which an agent's attention is directed toward the deep governing
structure of domain $d$ rather than surface-statistical regularities.

\subsection{Coupling to Schema Coherence}

The schema coherence ODE is updated to include attentional gating:
\begin{equation}\label{app:eq:sigma-ode}
\dot{\sigma}_A(d,t) = \rho \cdot P_A(d,t) \cdot \alpha_A(d,t) \cdot (1 - \sigma_A(d,t)) - \epsilon_{\sigma} \cdot \sigma_A(d,t) \cdot \Omega_{AI}(d,t)
\end{equation}

The multiplicative gating $\alpha_A \cdot P_A$ ensures that schema
coherence can only grow when attention is directed at principled
structure. This formalises the intuition that ``you can't learn what
you don't attend to.''

\subsection{Attentional Fidelity ODE}

\begin{equation}\label{app:eq:alpha-ode}
\dot{\alpha}_A(d,t) = \gamma \cdot C_A(d,t) \cdot (1 - \alpha_A(d,t)) - \zeta_{\alpha} \cdot \alpha_A(d,t) \cdot R_A^{\text{surface}}(d,t)
\end{equation}

where:
\begin{itemize}[nosep]
\item $C_A(d,t)$: Structural complexity of the training curriculum
\item $R_A^{\text{surface}}(d,t)$: Surface-reward pressure (e.g., from AI bypass)
\item $\gamma$: Attention growth rate constant
\item $\zeta_{\alpha}$: Surface-reward attention suppression coefficient
\end{itemize}

\subsection{Interpretation}

High $\alpha_A$ indicates that the agent's attention is consistently
directed at the compositional structure underlying the domain. Low
$\alpha_A$ indicates attention captured by surface features, statistical
regularities, or AI-provided shortcuts.

% ============================================================
\section{Extension 2: Collective Schema Field (Social Cognition)}\label{app:social}
% ============================================================

\subsection{Schema Legibility}

For agents $A$ and $B$, schema legibility quantifies how well $A$'s
schema coherence in domain $d$ can be understood by $B$:
\begin{equation}\label{app:eq:schema-legibility}
\mu_{AB}(d,t) = \sigma_A(d,t) \cdot \phi(d_A, d_B) \cdot \sigma_B(d_{\text{adj}},t) \in [0,1]
\end{equation}

where $\phi(d_A, d_B)$ is the domain similarity between $A$'s and $B$'s
representations.

\subsection{Collective Schema Field}

The collective schema field captures the bidirectional legibility
between two agents across two domains:
\begin{equation}\label{app:eq:collective-field}
\Sigma_{A,B}(d_1,d_2,t) = \mu_{AB}(d_1,t) \cdot \mu_{BA}(d_2,t) \cdot \phi(d_1,d_2)
\end{equation}

\subsection{Interpretation}

High $\Sigma_{A,B}$ indicates that agents $A$ and $B$ can effectively
communicate and align their schemas across domains $d_1$ and $d_2$.
This formalises the notion of ``shared understanding'' or ``common ground''
in social cognition.

% ============================================================
\section{Extension 3: Executive Control State $\Xi_A(t)$}\label{app:executive}
% ============================================================

\subsection{Definition}

The executive control state comprises three components:
\begin{equation}\label{app:eq:executive-state}
\Xi_A(t) = \{\Xi_A^P(t), \Xi_A^I(t), \Xi_A^F(t)\}
\end{equation}

\textbf{$\Xi_A^P$ --- Planning:} Degree to which the training trajectory
is consistent with the Σ-Model optimal multi-step arc.

\textbf{$\Xi_A^I$ --- Inhibition:} Probability of choosing the structural
route over AI bypass when both are available.

\textbf{$\Xi_A^F$ --- Cognitive Flexibility:} Degree to which the agent
detects and adapts to phase transition thresholds from internal evidence.

\subsection{Executive Control ODE}

\begin{equation}\label{app:eq:executive-ode}
\dot{\Xi}_A(t) = \kappa_P \cdot [P^*(t) - \Xi_A^P(t)] + \kappa_I \cdot [I^*(t) - \Xi_A^I(t)] + \kappa_F \cdot [F^*(t) - \Xi_A^F(t)]
\end{equation}

where $P^*(t)$, $I^*(t)$, $F^*(t)$ are target values for planning,
inhibition, and flexibility respectively.

\subsection{Interpretation}

Executive control modulates the agent's ability to: (a) plan multi-step
training trajectories, (b) inhibit the impulse to take AI shortcuts,
and (c) flexibly adapt when phase transitions are detected. Low
executive control predicts susceptibility to AI bypass and failure
to navigate phase transitions.

% ============================================================
\section{Extension 4: Self-Model of Schema Coherence $\hat{M}_A(d,t)$}\label{app:metacognition}
% ============================================================

\subsection{Definition}

The self-model $\hat{M}_A(d,t) \in [0,1]$ represents the agent's
estimate of its own schema coherence $\sigma_A(d,t)$.

\textbf{Calibration error:}
\begin{equation}\label{app:eq:calibration-error}
\zeta_A(d,t) = \hat{M}_A(d,t) - \sigma_A(d,t) \quad \text{(Calibration Error)}
\end{equation}

\subsection{Self-Model ODE}

\begin{equation}\label{app:eq:self-model-ode}
\dot{\hat{M}}_A(d,t) = \nu_M \cdot [\sigma_A(d,t) - \hat{M}_A(d,t)] - \xi_M \cdot \Omega_{AI}(d,t) \cdot \hat{M}_A(d,t)
\end{equation}

\subsection{Key Insight: AI Bypass Inflates Metacognitive Confidence}

The AI bypass risk $\Omega_{AI}$ appears in the self-model ODE with
a negative sign on $\hat{M}_A$, which means that when AI bypass is
active, the self-model is driven \emph{away} from true $\sigma_A$.
Agents using AI-provided outputs receive systematically inflated
performance feedback, producing overconfidence ($\zeta_A > 0$).

\subsection{Interpretation}

High calibration error $\zeta_A$ indicates metacognitive dysfunction:
the agent believes it has high schema coherence when it actually has
low schema coherence. This formalises the Dunning-Kruger effect in
the context of compositional generalisation.

% ============================================================
\section{Extension 5: Cross-Modal Schema Transfer}\label{app:multimodal}
% ============================================================

\subsection{The Domain $\times$ Modality Product Space}

All core state variables extend to the product space $D \times M$:
\begin{equation}\label{app:eq:modalities}
M = \{\text{text, visual, auditory, sensorimotor, symbolic}\}
\end{equation}

\subsection{Cross-Modal Schema Transfer Function}

\begin{equation}\label{app:eq:cross-modal-transfer}
\Theta_A(d, m_1, m_2, t) = \sigma_A(d, m_1, t) \cdot \omega(m_1, m_2)
\end{equation}

where $\omega(m_1, m_2)$ is the modality similarity weight.

\subsection{Extended Schema Coherence ODE}

\begin{equation}\label{app:eq:sigma-ode-v3}
\begin{aligned}
\dot{\sigma}_A(d,m,t) = {} & \rho \cdot P_A(d,m,t) \cdot \alpha_A(d,m,t) \cdot (1 - \sigma_A(d,m,t)) \\
& - \epsilon_{\sigma} \cdot \sigma_A(d,m,t) \cdot \Omega_{AI}(d,m,t) \\
& + \rho_{\theta} \cdot \sum_{m' \neq m} \Theta_A(d,m',m,t) \cdot \sigma_A(d,m',t)
\end{aligned}
\end{equation}

The cross-modal transfer term $\rho_{\theta} \cdot \sum \Theta_A \cdot \sigma_A$
enables schema coherence gained in one modality to facilitate schema
acquisition in another modality.

\subsection{Interpretation}

High $\sigma_A$ in one modality (e.g., text) facilitates transfer to
another modality (e.g., visual) proportional to the modality similarity
$\omega(m_1, m_2)$. This predicts that agents with high schema coherence
will show above-chance cross-modal transfer, while low-$\sigma_A$ agents
will not.

% ============================================================
\section{Extension 6: Benchmark Validity and Reliability}\label{app:validity}
% ============================================================

\subsection{Benchmark Validity Function}

\begin{equation}\label{app:eq:validity-function}
V_A(B,f,t) = CI(B,f) \cdot FD(B) \cdot DG(B) \cdot R_A(B,f,t)
\end{equation}

where:
\begin{itemize}[nosep]
\item $CI(B,f)$: Construct isolation --- degree to which benchmark $B$ isolates faculty $f$
\item $FD(B)$: Format diversity --- diversity of task formats in $B$
\item $DG(B)$: Difficulty gradient --- coverage of difficulty levels
\item $R_A(B,f,t)$: Reliability --- measurement stability across repeated assessments
\end{itemize}

\subsection{Benchmark Reliability Function}

\begin{equation}\label{app:eq:reliability}
R_A(B,f,t) = \begin{cases}
\displaystyle\frac{1}{1 + \frac{\text{Var}_k(\text{score}_A^k(B))}{E[\text{score}_A(B)]^2}} & \text{if } E[\text{score}_A(B)] > 0 \\[8pt]
0 & \text{if } E[\text{score}_A(B)] = 0
\end{cases}
\end{equation}

\subsection{Interpretation}

High $V_A$ indicates that benchmark $B$ is a valid and reliable measure
of faculty $f$ for agent $A$ at time $t$. Low $V_A$ indicates that the
benchmark may be confounded, narrow, or unreliable.

% ============================================================
\section{Summary of Extension Variables}\label{app:summary}
% ============================================================

\begin{center}
\begin{tabular}{c l l l}
\toprule
\textbf{Extension} & \textbf{Variable} & \textbf{Range} & \textbf{Role} \\
\midrule
Attention & $\alpha_A(d,t)$ & $[0,1]$ & Gates schema growth \\
Social Cognition & $\mu_{AB}(d,t)$ & $[0,1]$ & Schema legibility \\
Social Cognition & $\Sigma_{A,B}(d_1,d_2,t)$ & $[0,1]$ & Collective field \\
Executive Control & $\Xi_A^P(t)$ & $[0,1]$ & Planning quality \\
Executive Control & $\Xi_A^I(t)$ & $[0,1]$ & Inhibition quality \\
Executive Control & $\Xi_A^F(t)$ & $[0,1]$ & Cognitive flexibility \\
Metacognition & $\hat{M}_A(d,t)$ & $[0,1]$ & Self-model of $\sigma_A$ \\
Metacognition & $\zeta_A(d,t)$ & $[-1,1]$ & Calibration error \\
Multimodal & $\Theta_A(d,m_1,m_2,t)$ & $[0,1]$ & Cross-modal transfer \\
Validity & $V_A(B,f,t)$ & $[0,1]$ & Benchmark validity \\
Reliability & $R_A(B,f,t)$ & $[0,1]$ & Benchmark reliability \\
\bottomrule
\end{tabular}
\end{center}

% ============================================================
% Bibliography
% ============================================================
\bibliographystyle{plainnat}
\bibliography{../references}

\end{document}
